\chapter{Résolution exacte par une formulation compacte}
\section{Variables et objectif}
La formulation utilise le graphe non orienté complet $G=(V,E)$ et les arcs
orientés $D=\{(i,j):i\ne j\}$ associés aux arêtes. Elle s'appuie sur quatre
familles de variables :
\begin{itemize}
\item $z_i\in\{0,1\}$ indique si $i$ est une station ;
\item $y_{ij}\in\{0,1\}$ indique si le point $i$ est affecté à la station $j$ ;
\item $x_{ij}\in\{0,1\}$, pour $\{i,j\}\in E$, indique si cette arête appartient à l'anneau ;
\item $f_{ij}\geq0$, pour $(i,j)\in D$, désigne un flot fictif de connexité.
\end{itemize}
Le flot n'est pas un flux de voyageurs. C'est un outil mathématique qui empêche
la formation de plusieurs cycles disjoints. L'objectif est
\begin{equation}
\min\;\alpha\sum_{\{i,j\}\in E}d_{ij}x_{ij}
       +\sum_{i\in V}\sum_{j\in V}d_{ij}y_{ij}.
\end{equation}

\section{Sélection, affectation et degré}
Le nombre de stations et la présence du sommet imposé s'écrivent
\begin{equation}
\sum_{i\in V}z_i=p,\qquad z_r=1.
\end{equation}
Les affectations vérifient
\begin{align}
\sum_{j\in V}y_{ij}&=1 &&\forall i\in V,\\
y_{ij}&\leq z_j &&\forall i,j\in V,\\
y_{ii}&=z_i &&\forall i\in V.
\end{align}
Ces égalités imposent une affectation unique, exclusivement vers une station
ouverte, ainsi que l'auto-affectation des stations. L'objectif choisit une station
la plus proche à l'optimum, sans avoir à coder cette règle sous forme de
contraintes supplémentaires.

Chaque station possède deux arêtes d'anneau ; un point non retenu n'en possède
aucune :
\begin{equation}
\sum_{j\in V\setminus\{i\}}x_{\{i,j\}}=2z_i\qquad\forall i\in V.
\end{equation}
Le programme ajoute également les liaisons $x_{ij}\leq z_i$ et $x_{ij}\leq z_j$.
Dans une solution binaire, le degré exclut déjà les arêtes incidentes aux
non-stations, mais ces liaisons rendent cette relation explicite dans le modèle.

\section{Éliminer les cycles disjoints par un flot}
Les seules contraintes de degré permettraient plusieurs cycles séparés. Pour
obtenir un cycle unique, la racine émet $p-1$ unités de flot et chaque autre
station en consomme une :
\begin{align}
\sum_{j\ne r}f_{rj}-\sum_{j\ne r}f_{jr}&=p-1,\\
\sum_{j\ne i}f_{ji}-\sum_{j\ne i}f_{ij}&=z_i
   &&\forall i\ne r.
\end{align}
Le flot ne peut emprunter qu'une arête retenue :
\begin{equation}
f_{ij}+f_{ji}\leq(p-1)x_{ij}\qquad\forall\{i,j\}\in E.
\end{equation}
Les variables sont aussi bornées individuellement par $p-1$ dans le code.

Pour comprendre la connexité, supposons qu'un ensemble non vide de stations
forme une composante qui ne contient pas $r$. En additionnant ses équations de
conservation, les flots internes se compensent et la demande totale reste
strictement positive. Or aucune arête active ne relie cette composante au reste
du graphe, donc aucun flot ne peut y entrer. C'est une contradiction.
Le graphe actif est ainsi connexe. Comme ses sommets ont tous degré deux et
$p\geq3$, il forme un unique cycle.

Réciproquement, tout cycle réalisable admet un flot de $p-1,p-2,\ldots,1$
unités le long d'un sens du cycle, chaque station consommant une unité.

\section{Taille du modèle et outil de résolution}
Le modèle contient $n$ variables de station, $n^2$ d'affectation,
$n(n-1)/2$ d'arête et $n(n-1)$ de flot. Ses nombres de variables et de contraintes
sont d'ordre $n^2$ : c'est une formulation \emph{compacte}, ce qui ne rend pas
sa résolution entière polynomiale.

PuLP construit le modèle et appelle CBC, solveur de programmation entière
à séparation et évaluation arborescente\cite{cbc}. Le projet conserve le journal,
reconstruit le cycle et valide la solution. Il ne programme pas de séparation
de sous-tours ni d'arbre de Gomory-Hu ; la formulation non compacte reste
une extension possible.

\section{Une méthode exacte peut s'arrêter sans preuve}
Le programme contrôle les statuts du modèle et de la solution, vérifie
l'intégralité des décisions binaires puis recalcule le coût. Il annonce
\code{optimal} uniquement avec une preuve. Sous budget, une solution entière
sans preuve est classée \code{feasible_time_limited}. Sans solution utilisable,
aucun coût n'est inventé.

Si le journal fournit une borne inférieure $B$ compatible avec le coût
réalisable $C$, l'export conserve $g_B=(C-B)/\max(|C|,10^{-12})$.
Une borne absente reste absente. Cet écart se distingue de celui à un
optimum prouvé, défini au chapitre~\ref{chap:protocole}.
